% ============================================
% Introduction
% ============================================
\section{Introduction}
\label{sec:introduction}

In nature, species coexist and interact within complex communities, yet whether these assemblages function as cohesive, integrated units or merely as loose collections of independently acting species remains a fundamental question in ecology. Historically, this tension has been framed through two influential paradigms. Clements' ``superorganism'' view treats communities as discrete biological entities with emergent properties arising from species interdependencies, potentially as a result of coevolution\citep{Clements1916, Clements1936, Odum1969, Lovelock1979, WilsonSober1989}. In contrast, Gleason's individualistic hypothesis posits that species independently occupy niches, with community composition emerging from the coincidental overlap of species ranges\citep{Gleason1939, Cain1947, Mason1947, Whittaker1967}. Despite decades of research, the conditions that determine when communities behave as cohesive units versus loose species assemblages remain unclear.

These contrasting frameworks make distinct predictions in the context of community coalescence, the mixing of previously isolated communities\citep{Rillig2015, Lechon2021}. Coalescence occurs across diverse contexts and scales: environmental disturbances trigger microbial community mixing in soils\citep{Huet2023, Bresciani2025}, flooding events promote coalescence in aquatic and estuarine habitats\citep{Liu2024}, skin microbiomes undergo exchange through daily social contact\citep{Sarkar2024}, and gut microbiomes are subject to wholesale community transfer through fecal microbiota transplantation\citep{Xiao2020, Gupta2016}. Coalescence brings species with distinct interaction histories into contact, creating cross-community interactions that can reshape the resulting assemblage\citep{Rillig2015}. Thus, it provides a natural test for the individualistic versus holistic paradigms: The holistic view predicts that species within a community have correlated persistence, making the community the primary unit of selection, termed community-level selection. The individualistic view predicts species-level selection, yielding outcomes shaped by individual fitness regardless of community origin, with no systematic correlation in species persistence within parental communities.

Substantial theoretical and empirical evidence supports the holistic prediction that communities act as cohesive units during coalescence. Early theoretical work by Gilpin\citep{Gilpin1994} suggested that pre-assembled communities, having already undergone internal competitive exclusion, possess structured interaction patterns that differ fundamentally from randomly assembled communities. Because species within such communities have been filtered to coexist, their survival outcomes become coupled, producing asymmetric post-coalescence communities dominated by one parental community\citep{Debray2022}. Notably, this interaction structure can arise due to ecological exclusion alone, even in the absence of long-term coevolution, as demonstrated in synthetic microbial communities\citep{Venturelli2018, Maynard2020}. Building on this foundation, resource-consumer models formalized the mechanism by which such selection occurs, demonstrating that communities with superior collective fitness through resource consumption outcompete others in ways not predictable from individual species performance alone\citep{Tikhonov2016, Lechon2021, Xie2019, Xie2021}. Empirically, correlated selection between dominant and subdominant taxa has been observed across 100 coalescence experiments, confirming that species retention within communities is indeed correlated\citep{DiazColunga2022}. 

However, other evidence supports the individualistic prediction that species respond independently to coalescence events, with outcomes driven more by species-specific ecological responses to environmental conditions than by collective community dynamics. Empirical work reports that coexistence between species from different source communities via niche partitioning is more common than competitive exclusion during large-scale biogeographic interchange in marine and terrestrial biomes\citep{Vermeij1991}; similar patterns have been observed in microbial systems where closely related species coexist via resource partitioning\citep{Brochet2021}. In microbial systems, local environmental conditions can determine which species persist more strongly than exchange among neighboring communities, suggesting that environmental filtering of individual species can outweigh source-community identity\citep{VanderGucht2007}. Strain-resolved metagenomic analyses of \textit{in vitro} gut microbial communities showed species-level dynamics rather than community-level selection, with surviving species originating from both parental communities\citep{Goldman2025, Walton2025}. Together, these findings suggest that species-level selection may occur during coalescence when cross-community competitive exclusion is limited or when species-specific responses outweigh source-community identity in determining which species persist.

Together, these contrasting lines of evidence indicate that communities behave as units of selection in some coalescence events but not others; however, the factors governing these transitions remain poorly understood. Here, we combine empirical and theoretical approaches to investigate how effective interaction intensity and environmental feedbacks drive these contrasting outcomes. Using randomly assembled synthetic microbial communities across different interaction regimes, we classify post-coalescence outcomes into three types: Dominance (one community wins), Mixture (both persist), and Restructuring (a novel state emerges). Under weak interactions, Mixture prevails and species persistence is uncorrelated, consistent with individualistic dynamics. As effective interaction intensity increases, the system shifts toward Dominance, where species persistence becomes correlated, indicating community-level selection. A phenomenological generalized Lotka--Volterra (gLV) model with randomly sampled pairwise competitive interactions recapitulates the observed Dominance/Mixture transition, and a similar nutrient-dependent shift toward Dominance also appears in taxonomically richer communities derived from natural environmental samples. Further analysis reveals two mechanistic regimes underlying community-level selection: one where a few dominant taxa largely determine the winner, and another where collective multi-species dynamics shape the outcome. These results reconcile conflicting observations by identifying effective interaction intensity, environmental feedbacks, and assembly-generated interaction structure as key determinants of community-level selection.
